\section{Dijkstra's Algorithm Using a Set}
\label{sec:graph-dijkstra-set}

\problemstatement{Solve the same non-negative-weight shortest-path
problem as Section~\ref{sec:graph-dijkstra}, but use a \textbf{balanced
BST set} instead of a priority queue -- a variant that lets you actively
remove stale entries.}

\begin{patternbox}
The priority-queue Dijkstra tolerates stale duplicate entries and skips
them on pop. A \texttt{std::set} instead supports \textbf{erase}, so
when a node's distance improves we can \emph{delete} its old
$(\textit{dist},\textit{node})$ pair before inserting the new one --
emulating a decrease-key and keeping the structure free of outdated
copies.
\end{patternbox}

\subsection*{Set vs priority queue: the same algorithm, a different container}

The logic is identical to heap-based Dijkstra: repeatedly take the
smallest tentative-distance node, finalize it, and relax its edges. The
only difference is bookkeeping. A \texttt{set<pair<int,int>>} keeps
$(\textit{distance},\textit{node})$ pairs sorted, so
\texttt{begin()} is the minimum. When we find a shorter path to a
neighbor that already sits in the set with an older, larger distance, we
first \texttt{erase} that stale pair and then insert the improved one.
This keeps the set size bounded by $V$ rather than growing with every
improvement.

\begin{ideabox}[Erase-Then-Insert Emulates Decrease-Key]
The heap version cannot cheaply delete an interior element, so it lets
stale entries pile up and skips them later. The set version can delete
in $O(\log V)$, so it removes the obsolete pair up front. Both are
correct; the set trades a slightly larger constant factor for a smaller,
cleaner frontier.
\end{ideabox}

\subsection*{Algorithm}

\begin{enumerate}
  \item Initialize $\textit{dist}[\textit{src}]=0$, others $\infty$.
        Insert $(0,\textit{src})$ into the set.
  \item While the set is non-empty: take and erase the smallest pair
        $(\textit{d},\textit{node})$ from \texttt{begin()}.
  \item Relax each edge $\textit{node}\to\textit{nbr}$ of weight $w$: if
        $\textit{dist}[\textit{node}]+w<\textit{dist}[\textit{nbr}]$,
        erase the old $(\textit{dist}[\textit{nbr}],\textit{nbr})$ if
        present, update the distance, and insert the new pair.
\end{enumerate}

\subsection*{C++ implementation}

\begin{lstlisting}
#include <bits/stdc++.h>
using namespace std;

vector<int> dijkstraSet(int n, vector<vector<pair<int,int>>>& adj, int src) {
    vector<int> dist(n, INT_MAX);
    dist[src] = 0;
    set<pair<int,int>> st;               // {distance, node}, sorted
    st.insert({0, src});

    while (!st.empty()) {
        auto [d, node] = *st.begin();
        st.erase(st.begin());            // take the minimum

        for (auto& [nbr, w] : adj[node]) {
            if (dist[node] + w < dist[nbr]) {
                if (dist[nbr] != INT_MAX)
                    st.erase({dist[nbr], nbr});   // remove stale pair
                dist[nbr] = dist[node] + w;
                st.insert({dist[nbr], nbr});
            }
        }
    }
    return dist;
}

int main() {
    int n = 5;
    vector<vector<pair<int,int>>> adj(n);   // {neighbor, weight}
    auto addEdge = [&](int u, int v, int w) {
        adj[u].push_back({v, w});
        adj[v].push_back({u, w});
    };
    addEdge(0, 1, 4); addEdge(0, 2, 1); addEdge(2, 1, 2);
    addEdge(1, 3, 1); addEdge(2, 3, 5); addEdge(3, 4, 3);

    vector<int> dist = dijkstraSet(n, adj, 0);
    for (int d : dist) cout << d << " ";
    cout << endl;                        // 0 3 1 4 7
    return 0;
}
\end{lstlisting}

\subsection*{Dry run}

\dryrun{Same graph as Section~\ref{sec:graph-dijkstra}, source $0$.}

\begin{itemize}
  \item Set $\{(0,0)\}$. Take $(0,0)$; relax $1\!\to\!4$, $2\!\to\!1$.
        Set $\{(1,2),(4,1)\}$.
  \item Take $(1,2)$; relax $1$ via $2$: $3<4$, so erase $(4,1)$,
        insert $(3,1)$; relax $3\!\to\!6$. Set $\{(3,1),(6,3)\}$.
  \item Take $(3,1)$; relax $3$ via $1$: $4<6$, erase $(6,3)$, insert
        $(4,3)$. Set $\{(4,3)\}$.
  \item Take $(4,3)$; relax $4\!\to\!7$. Take $(7,4)$: nothing.
        Final $[0,3,1,4,7]$ -- no stale entries were ever skipped.
\end{itemize}

\subsection*{Complexity}

\begin{complexitybox}
\textbf{Time Complexity.} $O((V+E)\log V)$ -- same asymptotics as the
heap version; each edge may cause an erase and an insert, each
$O(\log V)$. \textbf{Space Complexity.} $O(V)$ -- the set never holds
more than one entry per node.
\end{complexitybox}

\begin{takeawaybox}
The set-based Dijkstra is the same greedy algorithm with a container
that supports deletion, so the frontier stays free of stale entries at
the cost of a larger constant. Use the priority-queue form by default;
reach for the set form when you specifically want an explicit
decrease-key or a frontier bounded by $V$.
\end{takeawaybox}
